\chapter{Resultados}
\label{chap:resultados}

Este capítulo sostiene que los resultados de un diagnóstico de cumplimiento solo tienen valor si cada cifra puede rastrearse hasta su evidencia y si lo que aún no se ha medido se declara como tal. Por ello, presenta los resultados en el orden en que se producen ---validación de los instrumentos, análisis documental, encuesta y brechas--- y para cada bloque muestra la tabla con la que se informará el resultado, el cálculo que la alimenta y su estado a la fecha de corte de este documento (7 de noviembre de 2026). El levantamiento se ejecuta entre noviembre de 2026 y febrero de 2027 (Sección~\ref{sec:cronograma}); hasta entonces, las celdas sin resultado llevan un guion largo (---) y ningún valor se estima ni se supone.

\section{Validación de los Instrumentos}
\label{sec:res_validacion_instrumentos}

La Tabla~\ref{tab:aiken_resultados} informa la validez de contenido del cuestionario y de la ficha documental, calculada con la V de Aiken (Ecuación~\ref{eq:aiken}) a partir de las calificaciones de los jueces (Anexo~\ref{anx:validacion_cuestionario}). Un ítem se conserva si $V \ge 0{,}80$ en los tres criterios.

\begin{table}[htbp]
  \centering
  \caption{Validez de contenido de los instrumentos (V de Aiken).}
  \label{tab:aiken_resultados}
  \scriptsize
  \renewcommand{\arraystretch}{1.2}
  \begin{tabular}{|l|c|c|c|c|c|}
    \hline
    \textbf{Bloque} & \textbf{Ítems} & \textbf{\shortstack{V mín.\\pertinencia}} & \textbf{\shortstack{V mín.\\claridad}} & \textbf{\shortstack{V mín.\\coherencia}} & \textbf{\shortstack{Ítems corregidos\\o eliminados}} \\
    \hline
    Ficha documental & 16 & --- & --- & --- & --- \\ \hline
    Acceso y datos generales (P00--A4) & 7 & --- & --- & --- & --- \\ \hline
    Preguntas de hechos (H1--H16) & 16 & --- & --- & --- & --- \\ \hline
    INF. Información recibida & 5 & --- & --- & --- & --- \\ \hline
    CON. Consentimiento & 5 & --- & --- & --- & --- \\ \hline
    DER. Conocimiento de derechos & 7 & --- & --- & --- & --- \\ \hline
    CNF. Confianza y seguridad & 6 & --- & --- & --- & --- \\ \hline
  \end{tabular}
  \par\vspace{2pt}
  \raggedright
  \textit{Nota:} --- indica un valor que se obtendrá del juicio de expertos; no disponible a la fecha de corte. Número de jueces: no disponible a la fecha de corte.
  \par\vspace{2pt}
  \elaboradopor{Ramón Alejandro Manzaba Castro}
  \fuente{Plantilla \path{encuesta/plantilla_juicio_expertos.xlsx}; \path{analisis/validacion.py}}
\end{table}

La Tabla~\ref{tab:ajustes_piloto} registra los ajustes que se deriven del juicio de expertos y de la prueba piloto, y la Tabla~\ref{tab:confiabilidad} informa la consistencia interna de cada escala (Ecuaciones~\ref{eq:cronbach}, \ref{eq:omega} y~\ref{eq:kr20}).

\begin{table}[htbp]
  \centering
  \caption{Ajustes del cuestionario tras el juicio de expertos y la prueba piloto.}
  \label{tab:ajustes_piloto}
  \footnotesize
  \renewcommand{\arraystretch}{1.2}
  \begin{tabularx}{\textwidth}{|c|>{\raggedright\arraybackslash}p{2.6cm}|L|L|}
    \hline
    \textbf{Ítem} & \textbf{Origen del ajuste} & \textbf{Redacción o decisión anterior} & \textbf{Redacción o decisión final} \\
    \hline
    --- & --- & --- & --- \\ \hline
  \end{tabularx}
  \par\vspace{2pt}
  \raggedright
  \textit{Nota:} el tiempo medio de respuesta del piloto y el número de participantes no están disponibles a la fecha de corte.
  \par\vspace{2pt}
  \elaboradopor{Ramón Alejandro Manzaba Castro}
  \fuente{Juicio de expertos y prueba piloto}
\end{table}

\begin{table}[htbp]
  \centering
  \caption{Confiabilidad de las escalas del cuestionario.}
  \label{tab:confiabilidad}
  \scriptsize
  \renewcommand{\arraystretch}{1.2}
  \begin{tabular}{|l|c|c|c|c|c|c|}
    \hline
    \textbf{Escala} & \textbf{Ítems} & \textbf{\shortstack{$\alpha$\\piloto}} & \textbf{\shortstack{$\omega$\\piloto}} & \textbf{\shortstack{$\alpha$\\muestra}} & \textbf{\shortstack{$\omega$\\muestra}} & \textbf{Valoración} \\
    \hline
    INF. Información recibida & 5 & --- & --- & --- & --- & --- \\ \hline
    CON. Consentimiento & 5 & --- & --- & --- & --- & --- \\ \hline
    DER. Conocimiento de derechos & 7 & --- & --- & --- & --- & --- \\ \hline
    CNF. Confianza y seguridad & 6 & --- & --- & --- & --- & --- \\ \hline
    Información recibida, hechos (H4--H7; KR-20) & 4 & \multicolumn{2}{c|}{---} & \multicolumn{2}{c|}{---} & --- \\ \hline
  \end{tabular}
  \par\vspace{2pt}
  \raggedright
  \textit{Nota:} valoración aceptable si el coeficiente es $\ge 0{,}70$ y buena si es $\ge 0{,}80$. --- indica un valor que se calculará con \path{analisis/validacion.py}; no disponible a la fecha de corte.
  \par\vspace{2pt}
  \elaboradopor{Ramón Alejandro Manzaba Castro}
  \fuente{Prueba piloto y encuesta; \path{analisis/validacion.py}}
\end{table}

Los resultados de validación son condición previa de todos los demás: si un ítem no supera el umbral de Aiken se corrige antes del piloto, y si una escala no alcanza la confiabilidad aceptable se informa ítem por ítem y no se resume en una puntuación de dimensión.

\section{Análisis Documental de las Universidades (OE2)}
\label{sec:res_documental}

La Tabla~\ref{tab:res_documental} informa, para cada indicador, cuántas de las 64 universidades y escuelas politécnicas cumplen (S), cumplen parcialmente (P) o no publican (N) el elemento verificado, junto con el cumplimiento documental $C_D$ de la Ecuación~\ref{eq:brecha}. Cada calificación está respaldada por la URL de la evidencia, la fecha de consulta y una copia en PDF registradas en la plantilla \path{analisis_documental/analisis_documental.xlsx}.

\begin{table}[htbp]
  \centering
  \caption{Resultados del análisis documental por indicador (64 universidades y escuelas politécnicas).}
  \label{tab:res_documental}
  \scriptsize
  \renewcommand{\arraystretch}{1.12}
  \begin{tabularx}{\textwidth}{|c|L|c|c|c|c|}
    \hline
    \textbf{ID} & \textbf{Indicador} & \textbf{S} & \textbf{P} & \textbf{N} & \textbf{$C_D$ (\%)} \\
    \hline
    ID01 & Política o aviso de privacidad publicado y accesible & --- & --- & --- & --- \\ \hline
    ID02 & Responsable identificado con dirección, teléfono y correo & --- & --- & --- & --- \\ \hline
    ID03 & Fines del tratamiento de datos de estudiantes & --- & --- & --- & --- \\ \hline
    ID04 & Base legal de cada tratamiento & --- & --- & --- & --- \\ \hline
    ID05 & Tiempo de conservación & --- & --- & --- & --- \\ \hline
    ID06 & Transferencias, destinatarios y garantías & --- & --- & --- & --- \\ \hline
    ID07 & Derechos del titular y forma de ejercerlos & --- & --- & --- & --- \\ \hline
    ID08 & Canal específico para ejercer derechos & --- & --- & --- & --- \\ \hline
    ID09 & Reclamo ante el responsable y la Autoridad & --- & --- & --- & --- \\ \hline
    ID10 & Revocación del consentimiento & --- & --- & --- & --- \\ \hline
    ID11 & Delegado de protección de datos con contacto & --- & --- & --- & --- \\ \hline
    ID12 & Referencia a la LOPDP y su Reglamento & --- & --- & --- & --- \\ \hline
    ID13 & Categorías especiales y su base de legitimación & --- & --- & --- & --- \\ \hline
    ID14 & Formularios con cláusula y consentimiento por finalidad & --- & --- & --- & --- \\ \hline
    ID15 & Aviso de \textit{cookies} con opción de rechazo (complementario) & --- & --- & --- & --- \\ \hline
    ID16 & Fecha de vigencia de la política (complementario) & --- & --- & --- & --- \\ \hline
  \end{tabularx}
  \par\vspace{2pt}
  \raggedright
  \textit{Nota:} --- indica un valor que se obtendrá del análisis documental (16 de noviembre a 24 de diciembre de 2026, con reverificación en enero de 2027); no disponible a la fecha de corte.
  \par\vspace{2pt}
  \elaboradopor{Ramón Alejandro Manzaba Castro}
  \fuente{Sitios web y documentos oficiales de las universidades; hoja Resumen de la plantilla documental}
\end{table}

La Tabla~\ref{tab:itd_sostenimiento} resume el índice de transparencia documental (Ecuación~\ref{eq:itd}) por tipo de sostenimiento, y la confiabilidad del registro se expresa con el porcentaje de acuerdo entre la primera verificación y la reverificación del 20\,\% de las universidades.

\begin{table}[htbp]
  \centering
  \caption{Índice de transparencia documental por tipo de sostenimiento.}
  \label{tab:itd_sostenimiento}
  \footnotesize
  \renewcommand{\arraystretch}{1.2}
  \begin{tabular}{|l|c|c|c|c|}
    \hline
    \textbf{Sostenimiento} & \textbf{Universidades} & \textbf{$ITD$ medio} & \textbf{$ITD$ mínimo} & \textbf{$ITD$ máximo} \\
    \hline
    Públicas & 36 & --- & --- & --- \\ \hline
    Particulares cofinanciadas & 8 & --- & --- & --- \\ \hline
    Particulares autofinanciadas & 20 & --- & --- & --- \\ \hline
    \textbf{Total} & \textbf{64} & --- & --- & --- \\ \hline
    \multicolumn{2}{|l|}{Acuerdo de la reverificación (13 universidades)} & \multicolumn{3}{c|}{---} \\ \hline
  \end{tabular}
  \par\vspace{2pt}
  \raggedright
  \textit{Nota:} el número de universidades por sostenimiento procede del portal del CES \cite{ces2026estadisticas}; --- indica un valor no disponible a la fecha de corte.
  \par\vspace{2pt}
  \elaboradopor{Ramón Alejandro Manzaba Castro}
  \fuente{Hojas Resumen y Reverificación de la plantilla documental}
\end{table}

La lectura de estos resultados se limita a lo que miden: qué informa cada universidad a cualquier titular en la fecha de consulta. Un resultado N no prueba que la universidad no trate los datos conforme a la ley, sino que el titular no puede comprobarlo con la información publicada, lo cual, para las obligaciones del artículo~12, equivale a un incumplimiento de la obligación de informar.

\section{Resultados de la Encuesta a Estudiantes (OE3)}
\label{sec:res_encuesta}

La Tabla~\ref{tab:muestra_obtenida} informa la participación obtenida frente a la meta de la Sección~\ref{sec:componente_encuesta}.

\begin{table}[htbp]
  \centering
  \caption{Participación en la encuesta y composición de la muestra.}
  \label{tab:muestra_obtenida}
  \footnotesize
  \renewcommand{\arraystretch}{1.2}
  \begin{tabular}{|l|c|c|}
    \hline
    \textbf{Concepto} & \textbf{Meta} & \textbf{Obtenido} \\
    \hline
    Respuestas recibidas & 423 & --- \\ \hline
    Excluidas por no aceptar, ser menor de edad o no ser estudiante & --- & --- \\ \hline
    Respuestas válidas & 384 & --- \\ \hline
    De universidades públicas & 246 (64\,\%) & --- \\ \hline
    De universidades privadas (cofinanciadas y autofinanciadas) & 138 (36\,\%) & --- \\ \hline
    No sabe el tipo de universidad & --- & --- \\ \hline
  \end{tabular}
  \par\vspace{2pt}
  \raggedright
  \textit{Nota:} recolección prevista del 7 de diciembre de 2026 al 15 de enero de 2027; --- indica un valor no disponible a la fecha de corte.
  \par\vspace{2pt}
  \elaboradopor{Ramón Alejandro Manzaba Castro}
  \fuente{Base anónima de la encuesta; \path{analisis/descriptivos.py}}
\end{table}

\subsection{Preguntas de hechos}

La Tabla~\ref{tab:res_hechos} informa los indicadores de hechos, en total y por tipo de universidad, con la prueba de asociación correspondiente (Ecuación~\ref{eq:chi}) corregida por comparaciones múltiples. La tabla se generará con \path{analisis/exportar_latex.py} a partir de los datos reales.

\begin{table}[htbp]
  \centering
  \caption{Indicadores de hechos de los estudiantes.}
  \label{tab:res_hechos}
  \scriptsize
  \setlength{\tabcolsep}{3pt}
  \renewcommand{\arraystretch}{1.12}
  \begin{tabularx}{\textwidth}{|c|L|c|c|c|c|c|}
    \hline
    \textbf{ID} & \textbf{Indicador (\% de estudiantes que\ldots)} & \textbf{Total} & \textbf{Pública} & \textbf{Privada} & \textbf{$p$ (Holm)} & \textbf{$V_C$} \\
    \hline
    IE01 & declaran que se les pidieron datos de salud, discapacidad o biométricos (H1) & --- & --- & --- & --- & --- \\ \hline
    IE02 & recibieron una explicación de la finalidad de esos datos (H9) & --- & --- & --- & --- & --- \\ \hline
    IE03 & fueron informados de la finalidad (H4) & --- & --- & --- & --- & --- \\ \hline
    IE04 & fueron informados del responsable y su contacto (H5) & --- & --- & --- & --- & --- \\ \hline
    IE05 & fueron informados del tiempo de conservación (H6) & --- & --- & --- & --- & --- \\ \hline
    IE06 & fueron informados de las comunicaciones a terceros (H7) & --- & --- & --- & --- & --- \\ \hline
    IE07 & declaran que se les pidió el consentimiento (H2) & --- & --- & --- & --- & --- \\ \hline
    IE08 & pudieron consentir por separado los usos no necesarios (H8) & --- & --- & --- & --- & --- \\ \hline
    IE11 & conocen el canal para ejercer sus derechos (H16) & --- & --- & --- & --- & --- \\ \hline
    IE12 & recibieron información o capacitación (H15) & --- & --- & --- & --- & --- \\ \hline
    IE14 & recibieron una comunicación de vulneración (H12) & --- & --- & --- & --- & --- \\ \hline
    IE15 & fueron informados de plataformas que elaboran perfiles (H10) & --- & --- & --- & --- & --- \\ \hline
  \end{tabularx}
  \par\vspace{2pt}
  \raggedright
  \textit{Nota:} --- indica un valor no disponible a la fecha de corte. Las pruebas de asociación son exploratorias (Sección~\ref{sec:plan_analisis}).
  \par\vspace{2pt}
  \elaboradopor{Ramón Alejandro Manzaba Castro}
  \fuente{Base anónima de la encuesta; \path{analisis/descriptivos.py} e \path{analisis/inferencial.py}}
\end{table}

\subsection{Dimensiones de Likert}

La Tabla~\ref{tab:res_dimensiones} resume las cuatro dimensiones con su mediana y rango intercuartílico, en total y por tipo de universidad, junto con la prueba de normalidad y la U de Mann-Whitney con su tamaño del efecto.

\begin{table}[htbp]
  \centering
  \caption{Dimensiones de percepción de los estudiantes.}
  \label{tab:res_dimensiones}
  \scriptsize
  \setlength{\tabcolsep}{3pt}
  \renewcommand{\arraystretch}{1.2}
  \begin{tabular}{|l|c|c|c|c|c|c|c|c|}
    \hline
    \textbf{Dimensión} & \textbf{Mdn} & \textbf{RIC} & \textbf{\shortstack{KS\\$p$}} & \textbf{\shortstack{Mdn\\pública}} & \textbf{\shortstack{Mdn\\privada}} & \textbf{$U$} & \textbf{\shortstack{$p$\\(Holm)}} & \textbf{$r$} \\
    \hline
    INF. Información recibida & --- & --- & --- & --- & --- & --- & --- & --- \\ \hline
    CON. Consentimiento & --- & --- & --- & --- & --- & --- & --- & --- \\ \hline
    DER. Conocimiento de derechos & --- & --- & --- & --- & --- & --- & --- & --- \\ \hline
    CNF. Confianza y seguridad & --- & --- & --- & --- & --- & --- & --- & --- \\ \hline
  \end{tabular}
  \par\vspace{2pt}
  \raggedright
  \textit{Nota:} Mdn = mediana; RIC = rango intercuartílico; KS = Kolmogorov-Smirnov con corrección de Lilliefors. --- indica un valor no disponible a la fecha de corte.
  \par\vspace{2pt}
  \elaboradopor{Ramón Alejandro Manzaba Castro}
  \fuente{Base anónima de la encuesta; \path{analisis/inferencial.py}}
\end{table}

\subsection{Contraste de hipótesis}

La Tabla~\ref{tab:res_hipotesis} informa el resultado de las tres hipótesis de la Sección~\ref{sec:hipotesis}, con $\alpha = 0{,}05$.

\begin{table}[htbp]
  \centering
  \caption{Contraste de las hipótesis H1 a H3.}
  \label{tab:res_hipotesis}
  \footnotesize
  \renewcommand{\arraystretch}{1.2}
  \begin{tabularx}{\textwidth}{|c|L|>{\raggedright\arraybackslash}p{2.6cm}|c|c|c|c|}
    \hline
    \textbf{Hip.} & \textbf{Enunciado} & \textbf{Prueba} & \textbf{Estadístico} & \textbf{$p$} & \textbf{Efecto} & \textbf{Decisión} \\
    \hline
    H1 & La información recibida (INF) difiere entre universidades públicas y privadas & U de Mann-Whitney & --- & --- & $r$ = --- & --- \\ \hline
    H2 & El conocimiento de derechos (DER) se relaciona con la confianza (CNF) & Spearman & --- & --- & $\rho$ = --- & --- \\ \hline
    H3 & La proporción a quienes se pidió consentimiento (H2) difiere según el tipo de universidad & Chi-cuadrado o Fisher & --- & --- & $V_C$ = --- & --- \\ \hline
  \end{tabularx}
  \par\vspace{2pt}
  \raggedright
  \textit{Nota:} --- indica un valor no disponible a la fecha de corte. La verificación cruzada en jamovi o JASP se registrará en \path{CAMBIOS.md}.
  \par\vspace{2pt}
  \elaboradopor{Ramón Alejandro Manzaba Castro}
  \fuente{Base anónima de la encuesta; \path{analisis/inferencial.py}}
\end{table}

Estos resultados describen a los estudiantes que respondieron y no a toda la población estudiantil. Su valor no está en estimar una proporción nacional, sino en mostrar, con el testimonio de los propios titulares, si las obligaciones de información, consentimiento y derechos llegan a quienes deben beneficiar, y si esa experiencia difiere según el tipo de universidad.

\section{Brechas entre la Norma, lo Publicado y lo Percibido (OE4)}
\label{sec:res_brechas}

La Tabla~\ref{tab:brechas} contrasta los siete pares de indicadores equivalentes con la Ecuación~\ref{eq:brecha}. La brecha normativa ($B_N$) mide cuánto falta para que todas las universidades publiquen el elemento exigido por la ley; la brecha de percepción ($B_P$) mide cuánto de lo publicado no llega a los estudiantes.

\begin{table}[htbp]
  \centering
  \caption{Brechas por par de indicadores equivalentes.}
  \label{tab:brechas}
  \footnotesize
  \renewcommand{\arraystretch}{1.2}
  \begin{tabular}{|l|c|c|c|c|c|}
    \hline
    \textbf{Elemento (par)} & \textbf{Norma} & \textbf{$C_D$ (\%)} & \textbf{$C_E$ (\%)} & \textbf{$B_N$} & \textbf{$B_P$} \\
    \hline
    Fines (ID03--H4) & 100 & --- & --- & --- & --- \\ \hline
    Responsable (ID02--H5) & 100 & --- & --- & --- & --- \\ \hline
    Conservación (ID05--H6) & 100 & --- & --- & --- & --- \\ \hline
    Terceros (ID06--H7) & 100 & --- & --- & --- & --- \\ \hline
    Consentimiento por finalidad (ID14--H8) & 100 & --- & --- & --- & --- \\ \hline
    Derechos (ID07--DER $\ge 4$) & 100 & --- & --- & --- & --- \\ \hline
    Canal de derechos (ID08--H16) & 100 & --- & --- & --- & --- \\ \hline
  \end{tabular}
  \par\vspace{2pt}
  \raggedright
  \textit{Nota:} la columna Norma expresa que la ley exige el elemento a todas las universidades y a todos los titulares \cite[arts.~8 y 12]{lopdp2021}; --- indica un valor no disponible a la fecha de corte.
  \par\vspace{2pt}
  \elaboradopor{Ramón Alejandro Manzaba Castro}
  \fuente{Tablas~\ref{tab:res_documental}, \ref{tab:res_hechos} y~\ref{tab:res_dimensiones}}
\end{table}

La interpretación de las brechas sigue tres reglas fijadas de antemano. Primera, una $B_N$ alta indica un incumplimiento de la obligación de informar en el plano documental. Segunda, una $B_P$ positiva y alta indica que la información existe pero no llega al titular ---por ejemplo, porque está en un documento difícil de encontrar o redactado sin claridad---, lo que compromete el principio de transparencia \cite[art.~10, lit.~c]{lopdp2021}. Tercera, una $B_P$ negativa indica que los estudiantes perciben más de lo que las universidades publican en su sitio web, lo que puede deberse a información entregada por otros medios (formularios impresos o comunicaciones directas) y se analiza con las respuestas de la dimensión INF del cuestionario. Las brechas más altas determinan qué escenarios de riesgo se priorizan en la metodología del Capítulo~\ref{chap:propuesta}, conforme al criterio de activación definido en \path{encuesta/plan_analisis.md}.

\section{Síntesis de Resultados}

A la fecha de corte, el capítulo permite afirmar con evidencia lo siguiente: los instrumentos y las fórmulas con los que se calcularán todos los resultados están definidos y verificados ---la plantilla documental contiene 1\,410 fórmulas que se evalúan sin errores y reproducen, con datos sintéticos de prueba, los cálculos independientes del índice, del cumplimiento por indicador y del acuerdo de la reverificación; los scripts de análisis en Python superan doce pruebas con valores de referencia y se ejecutaron de principio a fin con datos simulados, que no se usan como resultados---, y cada tabla de resultados está vinculada con su fuente, su ecuación y su indicador. Lo que el capítulo no afirma todavía es el grado de cumplimiento de las universidades ni la percepción de los estudiantes, porque ambos dependen del levantamiento programado entre noviembre de 2026 y febrero de 2027. Presentar esas cifras antes de medirlas contradiría el propósito mismo del trabajo, que es sustituir los supuestos por evidencia.
